\section{终局：跨宇宙、跨世界、跨本体}

本章讨论更远的终局模型。谢晨认为，最终模型应该 cross universe、cross world、cross embodiment，即跨宇宙、跨世界、跨本体。这个说法听起来宏大，但它背后有一个朴素目标：提升泛化性。人类可以从现实到游戏、从厨房到办公室、从手到工具快速迁移；机器人如果只能在一个仿真、一个房间或一个本体上工作，就还不是通用智能。

游戏数据在这里有特殊意义。游戏不一定物理真实，但它提供丰富世界、规则、目标和第一视角交互，适合预训练阶段提升 agent 跨世界能力。DeepMind 等机构大量使用游戏数据，正是因为跨宇宙训练可以让模型习惯在不同规则和观测分布中行动。它不能替代物理 Real2Sim，但能补足更广泛的 world diversity。

\begin{figure}[H]
\centering
\includegraphics[width=0.92\textwidth]{cross-universe-model.png}
\caption{跨宇宙、跨世界、跨本体：终局模型要提升跨环境泛化能力。自制概念图，依据 01:21:25--01:23:28 对谈内容整理。}
\end{figure}

\begin{knowledgebox}{读图：三种跨越分别解决不同泛化问题}
跨宇宙解决仿真、游戏和现实之间的 domain shift；跨世界解决家庭、工厂、酒店、餐馆等环境差异；跨本体解决不同机器人身体和动作空间差异。读图时要注意，三者不是口号，而是三类数据分布差异。终局模型要能在这些分布之间抽取任务本质。
\end{knowledgebox}

\begin{knowledgebox}{术语消化：三种跨越}
\begin{tabularx}{\textwidth}{p{0.18\textwidth}p{0.36\textwidth}X}
\toprule
能力 & 含义 & 训练要求 \\
\midrule
跨宇宙 & 在游戏、仿真和真实世界之间迁移 & 多样世界数据、规则学习和 domain adaptation。 \\
跨世界 & 在不同任务环境之间迁移 & 场景资产、任务分布和真实反馈要足够广。 \\
跨本体 & 在不同机器人身体之间迁移 & 动作空间、传感器、本体几何和控制约束要建模。 \\
\bottomrule
\end{tabularx}
\end{knowledgebox}

\subsection{还在 GPT-1 阶段}

本节回到现实阶段判断。谢晨认为具身智能整体仍在 GPT-1 阶段，意思不是完全没有能力，而是还没找到稳定 scaling law 配方。特斯拉 FSD 的类比很有帮助：早期不断加数据不一定持续提升，直到端到端路线打通，数据、算力和模型扩张才开始更稳定地带来能力提升。具身智能还没有出现这个明确时刻。

但他并不悲观。原因有三点：今天进入具身智能的创始团队和科学家密度远高于早期自动驾驶；资本和产业关注度更高；大模型和自动驾驶已经提供了 scaling、transformer、数据飞轮、RL 和端到端训练的经验。换句话说，当前是早期，但不是从零开始的早期。

\begin{figure}[H]
\centering
\includegraphics[width=0.96\textwidth]{embodied-gpt1-stage.png}
\caption{具身智能仍在 GPT-1 阶段：还没找到稳定 scaling law，数据 recipe 仍在探索。自制概念图，依据 01:23:28--01:28:21 对谈内容整理。}
\end{figure}

\begin{knowledgebox}{读图：GPT-1 阶段指的是 recipe 未收敛}
图中最重要的是“unknown recipe”。行业知道需要数据、算力、模型、仿真和 RL，但还不知道怎样的组合会稳定提升真实机器人能力。与其说 GPT-1 阶段意味着能力弱，不如说它意味着还没有形成 GPT-3/ChatGPT 那样清晰的规模化路径。
\end{knowledgebox}

\begin{importantbox}{Scaling law moment}
具身智能的 scaling law moment 不是融资、demo 或单点模型发布，而是团队发现：在一套架构和数据管线下，继续增加高质量数据、算力和环境，真实任务能力会可预测地提升。
\end{importantbox}

\subsection{Progress 比融资更健康}

本节保留访谈后半段的创业者判断。谢晨担心行业过分关注融资金额和估值，而忽略真实 progress。从 RL 角度看，如果团队的 reward model 变成“演讲、做 demo、吸引投资人”，行为就会被错误奖励牵引；更健康的 reward model 是服务客户、创造价值、收回付费，并用客户反馈持续改进产品。

这段话和技术主线其实一致：不管是数据、仿真还是公司，都需要真实反馈。对模型而言是真实部署反馈；对公司而言是客户付费和复购反馈；对行业而言是真实 progress，而不是 narrative progress。访谈中光轮从自动驾驶转向具身智能，也经历了从“我懂仿真”到“从具身本科重新学习”的过程，这说明高维问题不能直接套低维经验。

\begin{warningbox}{行业泡沫的一个信号}
如果团队主要优化融资叙事、demo 观感和估值，而不是客户价值、模型提升和真实部署，reward model 就偏了。对具身智能这样重工程行业，错误奖励会比短期技术失败更危险。
\end{warningbox}

\subsection{本章小结}

具身智能的终局是跨宇宙、跨世界和跨本体泛化；当前阶段仍在寻找 scaling law。行业也许会很快找到 recipe，但前提是把 reward 对准真实 progress：真实模型提升、真实客户价值和真实部署反馈。

